\documentclass[12pt]{article}

% Packages
\usepackage[
  paperwidth=18in,
  paperheight=5in,
  margin=0.5in
]{geometry}
\usepackage{amsmath, amssymb}
\usepackage{booktabs}
\usepackage{tabularx}

\begin{document}


% \clearpage
% \pdfpagewidth=18in
% \pdfpageheight=5in


\begin{table*}[h]
\centering
\small
\begin{tabular}{l c c c c}
\toprule
\textbf{Method} &
\textbf{Single / Multi Agent(LLM)} &
\textbf{What Is Optimized} &
\textbf{Prompt Exploration Strategy} &
\textbf{Information Provided in the Meta-Prompt used to generate the new prompt} \\
\midrule

\textbf{OPRO} \cite{OPRA}&
Single &
One prompt &
Iterative refinement &
\textbf{Parent prompts} together with their scalar performance scores \\

\textbf{EVOPROMPT} \cite{EVOPROMPT} &
Single &
One prompt &
Evolutionary search &
\textbf{Parent prompts only}; no scores or task feedback included \\

\textbf{PDO} \cite{PDO}&
Single &
One prompt &
Preference-based search &
\textbf{Parent prompts} with pairwise win/loss preference outcomes, plus fixed prompt mutation templates \\

\textbf{MIPRO} \cite{MIPRO}&
Single / Multi &
Multiple prompts &
Bayesian search &
\textbf{No parent prompts}; prompts are generated independently from task descriptions or templates \\

\textbf{GEPA} \cite{GEPA}&
Multi &
Multiple prompts &
Evolutionary search &
\textbf{Current system prompts} plus execution and evaluation traces from full system runs \\

\bottomrule
\end{tabular}
\caption{Comparison of prompt optimization methods across agent settings, optimization targets, and prompt exploration strategies. 
More details for each methods are given in Section~\ref{sec:summary}}
\label{tab:prompt_optimization_comparison}
\end{table*}

\section{Multi-agent system for Few-shot relation extraction with Optimized prompts}
\label{sec:method}

Our goal is the same as our previous work.  
We aim to improve few-shot relation extraction by adding more examples on top of a 1-shot example.  

We design a multi-agent system with three LLM agents.  
All agents use the same underlying LLM.  

The first agent is the \textbf{Example Generator}.  
This agent generates additional examples for a target relation.  

The second agent is the \textbf{Relation Inference Agent}.  
This agent performs relation inference as a binary classification task. 

The third agent is the \textbf{Evaluator (Critic)}.  
This agent evaluates the outputs of the Example Generator and the Relation Inference Agent.  

The Evaluator analyzes performance, inputs, and outputs of the other agents.  
It also reasons about success and failure cases.  

\textbf{The feedback from the Evaluator is used to update the prompts of the first and second agents.  
This process is repeated iteratively to improve prompt quality.}


% \clearpage
% \pdfpagewidth=8.5in
% \pdfpageheight=11in

\section{Summary of different Prompt Optimizer approaches}
\label{sec:summary}

\paragraph{OPRO:} (Simplest approach)
OPRO is a single-agent prompt optimization method that iteratively improves a single prompt.
It conditions prompt generation on a ranked history of previously evaluated prompts.
Each new prompt is generated by reasoning over parent prompts and their scalar performance scores.
The optimization proceeds through simple score-guided refinement without evolutionary operators.

\paragraph{EVOPROMPT:}
It is a single-agent evolutionary prompt optimization framework.
It maintains a population of prompts and evolves them using selection and mutation operators.
New prompts are generated solely from parent prompts, without access to scores or task feedback during generation.
Evaluation results are used only for population selection, not for prompt synthesis.

\paragraph{PDO:}
It optimizes a single prompt using preference-based feedback rather than explicit performance scores.
It relies on pairwise win–loss judgments produced by an LLM acting as a judge (They don't use ground-truth labels for evaluation of a prompt).
Prompt optimization is framed as a dueling bandit problem and solved using Double Thompson Sampling.
New prompt variants are generated using fixed, human-designed mutation templates applied to top-ranked prompts.

\paragraph{MIPRO}
MIPRO jointly optimizes multiple prompts across modules in a multi-stage LLM program.
It generates candidate instructions and demonstrations independently from task descriptions or templates.
No parent prompts or performance feedback are used during prompt generation.
A Bayesian surrogate model searches over combinations of candidates using only a task-level evaluation metric.

\paragraph{GEPA}
GEPA is a multi-agent evolutionary framework designed for compound AI systems with multiple LLM modules.  
It treats the entire set of system prompts as a single evolving population rather than optimizing prompts independently.  
In each iteration, a candidate system is selected from the population using Pareto-based selection.  
Within the selected system, a target module and its corresponding prompt are chosen in a round-robin manner.  
Prompt mutations are guided by natural-language reflection over execution and evaluation traces from full system runs.  
This design promotes diversity across modules and helps avoid premature convergence to local optima.


\bibliographystyle{plain}  
\bibliography{custom}    

\end{document}
